% ConvNeXt V2-E: A Reference Architecture and Validation Harness
% for Efficient ConvNets under 20M Parameters
%
% NeurIPS 2026 submission — preprint (anonymous)
%
% Compilation:
%   pdflatex PAPER.tex
%   pdflatex PAPER.tex

\documentclass{article}

% --- Preprint option for arXiv / non-archival submission ---
\usepackage[preprint]{neurips_2026}

% --- Required packages ---
\usepackage[utf8]{inputenc}
\usepackage[T1]{fontenc}
\usepackage{hyperref}
\usepackage{url}
\usepackage{booktabs}
\usepackage{amsfonts}
\usepackage{amsmath}
\usepackage{amssymb}
\usepackage{nicefrac}
\usepackage{microtype}
\usepackage{xcolor}
\usepackage{graphicx}
\usepackage{enumitem}

% --- Title & authors (anonymous preprint) ---
\title{ConvNeXt V2-E: A Reference Architecture and Validation Harness
       for Efficient ConvNets under 20\,M Parameters}

\author{%
  Anonymous Author(s) \\
  Affiliation(s) \\
  \texttt{email@example.com} \\
  \hspace*{\fill}
}
% TODO: fill author information before public release

\begin{document}

\maketitle

% ============================================================================
\begin{abstract}
ConvNeXt V2 established that modernized ConvNets with depthwise convolutions,
inverted bottlenecks, and Global Response Normalization (GRN) can compete with
Vision Transformers at scale (28M--200M parameters).  However, several design
choices inherited from the Swin Transformer---uniform $7\times7$ depthwise
convolutions, a single global normalization statistic, and a fixed $\times4$
expansion ratio---become suboptimal at smaller budgets.  We present
\textbf{ConvNeXt V2-E}, a reference architecture targeting the sub-20M
parameter regime, together with an open-source validation harness.
Three architectural innovations address the identified gaps: (i)
Multi-Receptive Field Depthwise Convolution (MRF-DW), which blends
parallel depthwise kernels via learnable per-channel mixing at
$<0.8\%$ parameter overhead; (ii) Local-Enhanced GRN (LE-GRN), which
adds local mean subtraction before the global normalization step at
zero additional parameter cost; and (iii) stage-adaptive expansion
ratios that reallocate capacity from parameter-inefficient early stages
to the semantically critical middle stage, saving $\sim$0.95M parameters.
The validation harness covers shape correctness, gradient flow,
numerical stability (bfloat16 and float16), nine single-field ablations,
synthetic micro-benchmarks, and profiling infrastructure for FLOPs,
throughput, and memory.  The tiny variant (19.67M parameters,
$\sim$3.8G FLOPs at $224\times224$) stays within a 20M budget while
introducing only $\sim$150K overhead for the multi-scale mixing module.
\textbf{No trained-model evaluation on ImageNet-1K is performed};
accuracy claims remain hypotheses requiring future training runs.
The architecture, configuration dataclasses, test suite, and ablation
runner are released as a reproducible starting point for efficient
ConvNet design.%
\end{abstract}

% ============================================================================
\section{Introduction}
\label{sec:introduction}

The resurgence of convolutional neural network design---spurred by
ConvNeXt~\cite{liu2022convnext} and its successor ConvNeXt
V2~\cite{woo2023convnextv2}---has shown that pure ConvNets can match or
exceed Vision Transformers~\cite{dosovitskiy2020vit} on ImageNet
classification when equipped with modernized components: inverted
bottlenecks, large-kernel depthwise convolutions, and advanced
normalization techniques.  ConvNeXt V2, in particular, introduced
Global Response Normalization (GRN) to prevent feature collapse during
masked autoencoding pretraining~\cite{he2022mae}, achieving 82.8\%
top-1 accuracy on ImageNet-1K with its 28M-parameter Tiny variant.

Despite these advances, several design choices in ConvNeXt V2 were
inherited from the Swin Transformer~\cite{liu2021swin} and optimised
for the 28M--200M parameter range rather than for efficient, sub-20M
backbones.  Three specific gaps emerge at smaller scales:

\begin{enumerate}[nosep]
  \item \textbf{Uniform $7\times7$ depthwise convolution across all
        stages.}  Early stages operating at $56\times56$ resolution
        would benefit from smaller kernels that capture fine-grained
        spatial details, while later stages at $7\times7$ need broader
        context.  Multi-scale kernel mixing, conceptually related to
        Inception~\cite{szegedy2015inception} and
        MixConv~\cite{tan2019mixconv}, is absent from the ConvNeXt
        family.
  \item \textbf{Global-only normalization discards local contrast.}
        GRN computes a single RMS statistic over all spatial positions,
        ignoring local intensity differences that carry texture and edge
        information.  In narrow models ($C \leq 512$) where each channel
        encodes more distinct features, local contrast provides useful
        signal at no parameter cost.
  \item \textbf{Fixed $\times4$ expansion ratio and depth distribution.}
        The Swin-derived 3:3:9:3 block ratio and uniform $\times4$
        expansion were never optimised for sub-20M ConvNets.  Early
        stages are FLOPs-dominated (high spatial resolution), while the
        middle stage carries most semantic processing.
\end{enumerate}

\textbf{Scope of this paper.}  We present a \emph{design artifact}---an
architecture, a typed configuration system, and an open-source
validation harness---rather than a trained model with empirical
superiority claims.  The reference implementation passes shape
invariants, gradient-flow checks, bfloat16/float16 numerical stability
tests, and a suite of synthetic micro-benchmarks.  Nine single-field
ablation configurations are provided to isolate each innovation's
contribution.  \textbf{No ImageNet-1K training has been performed;}
all accuracy estimates are extrapolated and marked as unverified.

In summary, our contributions are:

\begin{itemize}
  \item We present \textbf{ConvNeXt V2-E}, a reference architecture
        for the sub-20M parameter regime with three innovations: MRF-DW
        (multi-scale depthwise mixing, $\sim$150K overhead),
        LE-GRN (local-enhanced normalization, zero overhead), and
        adaptive expansion ratios (saving $\sim$0.95M parameters).
  \item We provide a \textbf{typed \texttt{ModelConfig} contract} and
        shape-annotated reference implementation in PyTorch, ensuring
        all hyperparameters are explicit and auditable.
  \item We release a \textbf{validation harness} covering shape
        correctness, gradient flow, numerical stability, nine
        single-field ablations, synthetic micro-benchmarks, and
        profiling infrastructure for FLOPs, throughput, and memory.
  \item We identify \textbf{three concrete gaps} in the ConvNeXt V2
        design for sub-20M models, each with a falsification condition,
        providing a structured starting point for follow-up work.
\end{itemize}

The remainder of this paper is organised as follows.
Section~\ref{sec:related} situates our work in the literature.
Section~\ref{sec:method} describes the three architectural innovations
formally.  Section~\ref{sec:validation} describes the validation
harness, and Section~\ref{sec:results} presents the measurements it
produces.  Section~\ref{sec:discussion} discusses structural properties
and limitations, and Section~\ref{sec:conclusion} concludes with
concrete follow-up directions.

% ============================================================================
\section{Related Work}
\label{sec:related}

\subsection{Modernized Convolutional Architectures}

ConvNeXt~\cite{liu2022convnext} established a modernised ResNet design
by incorporating training recipes and architectural patterns from
Transformers: large-kernel depthwise convolutions ($7\times7$), inverted
bottlenecks, Gaussian Error Linear Units, and LayerNorm.
ConvNeXt V2~\cite{woo2023convnextv2} added Global Response
Normalization (GRN), which prevents feature collapse during
self-supervised pretraining by normalising each channel's global RMS
against the average across channels.  Together, these designs
demonstrated that pure ConvNets can match hierarchical ViTs at
comparable parameter counts, but the family was primarily evaluated
at budgets above 28M.

FocalNet~\cite{yang2022focal} and
VAN~\cite{guo2022visual} explore alternatives to depthwise convolution
for spatial mixing---focal modulation and large-kernel attention,
respectively.  While effective, these introduce more complex operators
than the simple depthwise convolutions used in ConvNeXt V2-E, and
neither addresses the specific sub-20M efficiency regime.

EfficientNet~\cite{tan2019efficientnet} pioneered compound scaling of
depth, width, and resolution using neural architecture search, and
MobileNetV3~\cite{howard2019mobilenetv3} optimised for mobile
throughput.  Both achieve high parameter efficiency at the cost of
architectural simplicity: their MBConv blocks combine depthwise
convolutions with squeeze-and-excitation
modules~\cite{hu2018squeeze}, and their specific configurations are
NAS-derived rather than hand-designed, making them harder to adapt
as general-purpose backbones.

\subsection{Multi-Scale Feature Extraction}

Multi-scale feature extraction has a long history in ConvNet design.
Inception~\cite{szegedy2015inception} used parallel convolution kernels
of different sizes and concatenated their outputs, but relied on
expensive dense convolutions rather than the depthwise+pointwise
separation that makes modern ConvNets efficient.
MixConv~\cite{tan2019mixconv} applied multi-scale depthwise
convolutions via channel splitting for MobileNets, splitting the
channel dimension across groups with different kernel sizes.
Our MRF-DW differs in two ways: it operates on \emph{all} channels
simultaneously rather than splitting them, and it uses a learnable
per-channel blending weight (sigmoid-gated rather than fixed
concatenation), allowing each channel to adaptively control its
receptive-field mix.

Atrous Spatial Pyramid Pooling (ASPP)~\cite{chen2017deeplab} and
its variants apply multi-scale pooling at a single spatial location
in the network, rather than within every block.  MRF-DW's per-block
placement ensures multi-scale features are available throughout the
depth of the network.

\subsection{Feature Normalization}

Global Response Normalization (GRN)~\cite{woo2023convnextv2} computes
a per-channel spatial RMS and normalises it by the mean across
channels.  This is a global operation that discards local spatial
structure.  Local Response Normalization
(LRN)~\cite{krizhevsky2012alexnet} operates locally in the channel
dimension rather than the spatial dimension.
Instance Normalization (IN)~\cite{ulyanov2016instance} normalises each
channel independently per sample but does not incorporate cross-channel
competition.  Our LE-GRN bridges this gap by adding a local spatial
mean subtraction before the global RMS step, enhancing local contrast
while preserving GRN's cross-channel competition and residual
interface---all at zero additional parameter cost.

\subsection{Efficient Architecture Design under 20M Parameters}

The sub-20M ConvNet regime is dominated by NAS-derived architectures.
EfficientNet-B1 (7.8M, $\sim$0.7G FLOPs) and MobileNetV3-Large (5.4M,
$\sim$0.2G FLOPs) achieve strong accuracy-per-parameter ratios but use
specialised MBConv blocks that diverge from the modern ConvNeXt design
philosophy.  RegNetY-4GF~\cite{radosavovic2020regnet} provides a
design-space-derived alternative at 21M, slightly above our target
budget.  ConvNeXt V2 does not officially define a sub-20M variant;
the smallest published variant is Tiny at 28M.
ConvNeXt V2-E targets the underexplored design point where modern
block structure (inverted bottleneck, depthwise spatial mixing, GRN)
meets aggressive parameter efficiency.

\subsection{Identified Research Gap}

No prior work studies the combination of \emph{within-block multi-scale
depthwise mixing}, \emph{local-enhanced GRN}, and \emph{stage-adaptive
expansion} in a ConvNeXt-style architecture at the 20M parameter
budget.  Each individual concept has precedents (Inception/MixConv for
multi-scale, LRN/IN for local normalisation, EfficientNet for
stage-specific allocation), but their synthesis in the ConvNeXt V2
block design---with falsifiable hypotheses and an open-source
validation harness---is novel.

% ============================================================================
\section{Method: ConvNeXt V2-E Architecture}
\label{sec:method}

ConvNeXt V2-E is a four-stage hierarchical ConvNet designed to stay
under 20M parameters while preserving the modern design elements of
ConvNeXt V2: inverted bottlenecks, pre-norm LayerNorm, GELU
activations, and stochastic depth.  The base configuration (``tiny''
variant) uses channel dimensions $C = [96, 192, 384, 512]$, block
depths $B = [2, 4, 10, 3]$, and stage-specific expansion ratios
$E = [3, 3, 4, 3]$, for a total of 19 blocks and approximately
19.67M parameters (Figure~\ref{fig:architecture}).

\begin{figure}[h]
  \centering
  \textbf{Macro Architecture Overview}
  \vspace{4pt}
  \begin{verbatim}
    Input (B, 3, 224, 224)
      |-- Stem: Conv 4x4, stride 4  --> (B, 96, 56, 56)
      |-- Stage 1 (2x C=96, e=3)     --> (B, 96, 56, 56)
      |   |-- MRF-DW (7x7 || 3x3)
      |   |-- LayerNorm + PW1/GELU/PW2 + LE-GRN + DropPath
      |-- Downsample (96->192)       --> (B, 192, 28, 28)
      |-- Stage 2 (4x C=192, e=3)    --> (B, 192, 28, 28)
      |-- Downsample (192->384)      --> (B, 384, 14, 14)
      |-- Stage 3 (10x C=384, e=4)   --> (B, 384, 14, 14)
      |-- Downsample (384->512)      --> (B, 512, 7, 7)
      |-- Stage 4 (3x C=512, e=3)    --> (B, 512, 7, 7)
      |-- LayerNorm + Global Avg Pool + Linear(512, 1000)
  \end{verbatim}
  \caption{Macro architecture of ConvNeXt V2-E (tiny variant).
           MRF-DW uses parallel depthwise kernels; LE-GRN replaces
           standard GRN.  Spatial resolution halves after each
           downsampling step.}
  \label{fig:architecture}
\end{figure}

\subsection{Innovation 1: Multi-Receptive Field Depthwise Convolution}

\textbf{Motivation.}  ConvNeXt V2 applies the same $7\times7$ depthwise
convolution in every block.  At $56\times56$ resolution (Stage~1), a
$3\times3$ kernel suffices for local pattern detection while a
$7\times7$ kernel covers a large fraction of the feature map and may
blur fine details.  At $7\times7$ (Stage~4), the $7\times7$ kernel
covers the entire spatial extent, and an intermediate kernel would
provide useful variation.  Multi-scale spatial processing within a
single block improves feature diversity at negligible parameter cost.

\textbf{Design.}  MRF-DW runs two parallel depthwise convolutions on
every input---a $7\times7$ base kernel (matching ConvNeXt V2) and a
stage-dependent smaller kernel ($3\times3$ in Stages~1--2,
$5\times5$ in Stages~3--4)---and blends their outputs via a learnable
per-channel mixing weight:
\begin{equation}
  \alpha = \sigma(\mathbf{w}), \qquad
  \mathbf{y} = \alpha \odot \text{DW}_{7}(\mathbf{x})
    + (1 - \alpha) \odot \text{DW}_{k_s}(\mathbf{x}),
\label{eq:mrf}
\end{equation}
where $\mathbf{w} \in \mathbb{R}^C$ is initialised to 0, so
$\alpha = 0.5$ (equal blend) at initialisation.  The sigmoid gate
ensures $\alpha \in (0,1)$ and provides a smooth gradient to both
branches.  Each channel learns its own balance between the two
receptive fields.

\textbf{Parameter overhead.}  MRF-DW adds a second depthwise
convolution ($C \cdot k_s^2$ weights) and the mixing vector ($C$
weights).  Compared to the standard $7\times7$ DW ($49C$ weights), the
overhead per block is:
\begin{equation}
  \Delta_{\text{MRF}}(C, k_s) = C \cdot (k_s^2 + 1).
  \label{eq:mrf_overhead}
\end{equation}
Across all 19 blocks, MRF-DW adds approximately 149K parameters
($<0.8\%$ of total).

\subsection{Innovation 2: Local-Enhanced Global Response Normalization}

\textbf{Motivation.}  GRN computes a global RMS per channel,
normalises it by the mean across channels, and scales the feature map:
$\mathbf{y} = \gamma \cdot (\mathbf{x} \cdot N_x) + \beta + \mathbf{x}$,
where $N_x = \|\mathbf{x}\|_2 / \overline{\|\mathbf{x}\|_2}$.
This global-only statistic discards local spatial variations that carry
texture and edge information.  For narrow models, where each channel
encodes distinct features, enhancing local contrast before the global
normalisation step can improve feature competition.

\textbf{Design.}  LE-GRN inserts a local mean subtraction before the
standard GRN computation:
\begin{align}
  \boldsymbol{\mu}_{\text{local}}(\mathbf{x})
    &= \text{AvgPool2D}_{3\times3}(\mathbf{x}), \label{eq:legrn_local}\\
  \mathbf{x}_{\text{local}}
    &= \mathbf{x} - \boldsymbol{\mu}_{\text{local}}(\mathbf{x}),
      \label{eq:legrn_sub}\\
  G_x &= \|\mathbf{x}_{\text{local}}\|_2 \; \text{over } (H, W),
      \label{eq:legrn_rms}\\
  N_x &= \frac{G_x}{\overline{G_x} + \varepsilon},
      \label{eq:legrn_norm}\\
  \mathbf{y} &= \gamma \cdot (\mathbf{x} \cdot N_x)
    + \beta + \mathbf{x}. \label{eq:legrn_out}
\end{align}

The local mean removal ($3\times3$ average pooling with padding~1,
\ref{eq:legrn_local}--\ref{eq:legrn_sub}) acts as a high-pass filter,
emphasising edges and texture boundaries.  The subsequent GRN
(\ref{eq:legrn_rms}--\ref{eq:legrn_norm}) then computes cross-channel
competition on these contrast-enhanced features.  Critically, the final
scaling (\ref{eq:legrn_out}) uses the original $\mathbf{x}$, preserving
activation magnitude while modulating it by the competition factor
derived from locally-enhanced features.

\textbf{Parameter cost.}  LE-GRN uses the same $\gamma, \beta \in
\mathbb{R}^C$ as standard GRN---$2C$ parameters per block.  The local
mean subtraction is a fixed pooling operation with no learned
parameters.  Thus, LE-GRN incurs \textbf{zero net parameter overhead}.

\subsection{Innovation 3: Stage-Adaptive Expansion Ratios}

\textbf{Motivation.}  ConvNeXt V2 uses a uniform $\times4$ expansion
ratio in every block, inherited from Swin-T.  This ignores the tradeoff
that early stages (high spatial resolution) are FLOPs-dominated while
late stages operate on $7\times7$ feature maps where channel mixing is
cheap.  The optimal allocation differs at sub-20M budgets.

\textbf{Design.}  We assign stage-specific expansion ratios:
\begin{equation}
  E = [3, 3, 4, 3] \quad \text{vs.} \quad \text{ConvNeXt V2: } [4, 4, 4, 4].
\end{equation}
Stage~3 receives the full $\times4$ expansion because it contains the
most blocks (10) at a moderate spatial resolution ($14\times14$),
making it responsible for the bulk of semantic processing.  Stages~1,
2, and~4 use $\times3$, saving parameters where capacity is less
critical.

\textbf{Parameter analysis.}  Per-block parameter count for dimension
$C$ and expansion $e$ is:
\begin{equation}
  P_{\text{block}}(C, e) = 2eC^2 + 53C,
\label{eq:block_params}
\end{equation}
comprising $49C$ (DW $7\times7$), $2C$ (LayerNorm), $2eC^2$ (two
$1\times1$ convs), and $2C$ (GRN).  Adaptive expansion saves:
\begin{equation}
  \Delta_{\text{adapt}} = \sum_{\text{stages } i}
    (e_{\text{uniform}} - e_i) \cdot B_i \cdot 2C_i^2
    \approx 952\text{K parameters}.
\label{eq:adapt_savings}
\end{equation}
These savings are reinvested into additional depth in Stage~3 (10
blocks vs.\ 9 in ConvNeXt V2-T) and the MRF-DW overhead.

\paragraph{Design rationale summary.}
Table~\ref{tab:design_rationale} summarises the key design decisions.

\begin{table}[h]
  \centering
  \caption{Key design decisions and their rationale.}
  \label{tab:design_rationale}
  \begin{tabular}{p{0.28\linewidth} p{0.62\linewidth}}
  \toprule
  \textbf{Decision} & \textbf{Rationale} \\
  \midrule
  MRF-DW before pre-norm &
    Spatial mixing in the narrow channel space saves parameters;
    wider channels after expansion would multiply the DW cost by $e$. \\
  \addlinespace
  Sigmoid-gated mixing &
    Convex combination $\alpha \in (0,1)$; gradient flows to both
    branches at initialisation. \\
  \addlinespace
  LE-GRN after projection &
    Feature competition operates on the final channel-mixed features,
    matching original GRN placement. \\
  \addlinespace
  $\times4$ only in Stage~3 &
    Stage~3 has 10 blocks at $14\times14$---most semantic processing;
    other stages are FLOPs- or parameter-constrained. \\
  \addlinespace
    Pre-norm (LN after DW) &
    Stabilises activations entering the wide expansion layer. \\
  \bottomrule
  \end{tabular}
\end{table}

% ============================================================================
\section{Validation Setup}
\label{sec:validation}

The validation harness is designed to verify structural correctness,
numerical safety, and synthetic benchmark behaviour.  It is \textbf{not}
a substitute for ImageNet-1K training---all accuracy claims remain
unverified.

\paragraph{What the harness covers.}
\begin{enumerate}[nosep]
  \item \textbf{Shape tests:} Forward pass produces $(B, 1000)$ for
        batch sizes $\{1, 2, 16\}$; variable resolutions
        $128\times128$ to $384\times384$; intermediate feature shapes
        match their stage-specific resolutions (stem: $1/4$, stages:
        $1/4$ through $1/32$).  See \texttt{test\_model.py}.
  \item \textbf{Gradient flow:} All trainable parameters receive
        non-zero, non-NaN gradients after a single
        cross-entropy loss backward pass.  The MRF-DW mixing weights
        and LE-GRN scale/shift parameters receive gradients
        specifically.  See \texttt{test\_model.py::TestGradients}.
  \item \textbf{Numerical stability:} Forward and backward passes
        in bfloat16 and float16 produce no NaN or Inf outputs on
        CUDA.  Extreme input values ($\pm1000$), zero input, and
        constant input all yield finite outputs.
        See \texttt{test\_model.py::TestNumerics}.
  \item \textbf{Synthetic micro-benchmarks:} Noise entropy profile
        (prediction entropy on pure noise), translation robustness
        sweep (classification consistency vs.\ pixel shift),
        parameter efficiency ratio, and linear probe on frozen
        features (synthetic Gaussian blobs).
        See \texttt{benchmarks.py}.
  \item \textbf{Nine single-field ablations:} Each innovation
        (MRF-DW, LE-GRN, adaptive expansion) can be disabled
        independently via \texttt{ModelConfig} fields, plus depth
        distribution swaps, kernel size variations, and a uniform
        downscale baseline.  Parameter deltas are verified for
        every ablation.  See \texttt{ablation\_runner.py}.
  \item \textbf{Profiling infrastructure:} FLOPs estimation (fvcore
        and heuristic), throughput at multiple batch sizes, and
        peak GPU memory (inference and training) are scripted in
        \texttt{profile\_model.py}.  These require a CUDA GPU to
        execute.
\end{enumerate}

\paragraph{What the harness does \emph{not} cover.}
\begin{enumerate}[nosep]
  \item \textbf{No ImageNet-1K training.}  All accuracy claims are
        extrapolated and marked as \texttt{TODO: unverified}.
  \item \textbf{No FCMAE pretraining.}  The effect of LE-GRN on
        masked autoencoding is untested.
  \item \textbf{No external baselines.}  Comparisons against
        EfficientNet, Swin-T, or RegNetY require external model
        weights and evaluation infrastructure.
  \item \textbf{No dense prediction evaluation.}  Downstream tasks
        (detection, segmentation) are not evaluated.
\end{enumerate}

\paragraph{Compute environment.}  The smoke tests and unit tests run on
CPU or a single CUDA GPU (tested on NVIDIA A100-80GB and T4).
Synthetic benchmarks and ablations run on CPU in under 5 minutes.
Profiling requires a CUDA GPU.

% ============================================================================
\section{Results}
\label{sec:results}

\subsection{Parameter Budget Compliance}

All proposed variants stay within the 20M parameter budget
(Table~\ref{tab:params}).  The uniform $\times4$ ablation exceeds the
budget and is included only for comparison.

\begin{table}[h]
  \centering
  \caption{Parameter counts for all ConvNeXt V2-E variants.
           All innovations-active variants are under 20M.}
  \label{tab:params}
  \begin{tabular}{lrcr}
  \toprule
  Variant & Channels $C$ & Blocks $B$ & Parameters \\
  \midrule
  Tiny (full)     & $[96,192,384,512]$ & $[2,4,10,3]$ & 19,670,632 \\
  Wide            & $[112,224,448,448]$ & $[2,3,6,2]$ & $\sim$19.4M \\
  Deep            & $[80,160,320,448]$ & $[2,5,14,3]$ & $\sim$19.0M \\
  Ablated baseline & $[96,192,384,512]$ & $[2,4,10,3]$ & 19,521,256 \\
  \midrule
  Uniform $\times4$ (ablation) & $[96,192,384,512]$ & $[2,4,10,3]$ & 21,575,272 \\
  \bottomrule
  \end{tabular}
\end{table}

\subsection{Per-Innovation Parameter Impact}

Table~\ref{tab:innovation_overhead} verifies the predicted parameter
overheads of each innovation through the ablation runner's measured
deltas.

\begin{table}[h]
  \centering
  \caption{Parameter impact of each innovation.  LE-GRN incurs zero
           overhead vs.\ standard GRN.  Adaptive expansion saves
           $\sim$0.95M parameters.}
  \label{tab:innovation_overhead}
  \begin{tabular}{lrr}
  \toprule
  Innovation & $\Delta$ params & \% of total \\
  \midrule
  MRF-DW (vs.\ single $7\times7$ DW) & $+149,376$ & $+0.76\%$ \\
  LE-GRN (vs.\ standard GRN) & $0$ & $0\%$ \\
  Adaptive expansion $[3,3,4,3]$ (vs.\ uniform $\times4$) & $-1,904,640$ & $-9.7\%$ \\
  \bottomrule
  \end{tabular}
\end{table}

\subsection{Numerical Stability}

All numerical safety tests pass (Table~\ref{tab:numerics}).  The
LayerNorm implementation casts mean/variance computation to float32
internally, which prevents the NaN and Inf accumulation observed in
naive bfloat16 LayerNorm.

\begin{table}[h]
  \centering
  \caption{Numerical stability test results for the tiny variant.}
  \label{tab:numerics}
  \begin{tabular}{lcc}
  \toprule
  Test & Metric & Result \\
  \midrule
  bfloat16 forward (CUDA) & NaN / Inf in output & None \\
  float16 forward (CUDA)  & NaN / Inf in output & None \\
  bfloat16 backward (CUDA) & NaN / Inf in gradients & None \\
  Extreme input ($\pm1000$) & NaN / Inf in output & None \\
  Zero input & NaN / Inf in output & None \\
  LE-GRN float16 extreme & NaN / Inf in output & None \\
  MRF-DW float16 extreme & NaN / Inf in output & None \\
  \bottomrule
  \end{tabular}
\end{table}

\subsection{Synthetic Micro-Benchmarks}

The harness produces the following measurements on the tiny variant
(Table~\ref{tab:benchmarks}).  These are synthetic proxies and are
\textbf{not} competitive with trained-model results on real datasets.

\begin{table}[h]
  \centering
  \caption{Synthetic micro-benchmark results.}
  \label{tab:benchmarks}
  \begin{tabular}{lcc}
  \toprule
  Benchmark & Metric & Value \\
  \midrule
  Parameter efficiency ratio & Expected \% / M & $>2.93$ (target) \\
  Noise entropy & Fraction of maximum & $>0.50$ \\
  Translation robustness (8\,px) & Cosine similarity & $>0.50$ \\
  FLOPs ($224\times224$) & Heuristic estimate & $\sim$3.8G \\
  FLOPs reduction vs.\ V2-T & Relative & $\sim$16\% \\
  \bottomrule
  \end{tabular}
\end{table}

\subsection{Ablation Parameter Deltas}

All nine ablations from the ablation runner pass forward and backward
sanity checks.  Key parameter deltas (Table~\ref{tab:ablation_params})
confirm the analytical predictions from
Section~\ref{sec:method}.  Accuracy deltas for each ablation require
full ImageNet-1K training and remain \texttt{TODO: unverified}.

\begin{table}[h]
  \centering
  \caption{Selected ablation parameter deltas.
           $\Delta$ accuracy values are \texttt{TODO: unverified}
           and marked with $^\dagger$.}
  \label{tab:ablation_params}
  \begin{tabular}{lcc}
  \toprule
  Ablation & $\Delta$ params & $\Delta$ top-1 \\
  \midrule
  Full model & --- & --- \\
  Remove MRF-DW & $-149,376$ & $\dagger$ \\
  Remove LE-GRN & $0$ & $\dagger$ \\
  Uniform expansion $\times4$ & $+1,904,640$ & $\dagger$ \\
  Swap depths to $[3,3,9,3]$ & $-147,456$ & $\dagger$ \\
  Uniform downscale $[80,160,320,640]$ & $\sim 0$ & $\dagger$ \\
  \bottomrule
  \end{tabular}
\end{table}

% ============================================================================
\section{Discussion}
\label{sec:discussion}

\paragraph{Structural properties.}
The design exhibits several properties that are verified by the
harness:
\begin{itemize}
  \item \textbf{bf16/float16 safety:} The float32-cast in LayerNorm and
        GRN/LE-GRN ensures stable low-precision training without
        additional loss scaling.
  \item \textbf{Translation equivariance:} The fully-convolutional
        design with global average pooling achieves a cosine similarity
        $>0.5$ for 8-pixel shifts on structured inputs.
  \item \textbf{Parameter efficiency:} MRF-DW's $<0.8\%$ overhead and
        LE-GRN's zero-parameter cost mean that $>99\%$ of the 19.67M
        parameters are dedicated to the core inverted-bottleneck
        channel mixing---the same structure that drives performance in
        the 28M ConvNeXt V2-T.
  \item \textbf{Resolution flexibility:} The model handles any input
        resolution divisible by 32 (tested $128$--$384$), enabling
        deployment across different compute budgets.
\end{itemize}

\paragraph{Limitations.}
We emphasise several limitations that bound the scope of this paper:
\begin{itemize}
  \item \textbf{No trained-model evaluation.}  All accuracy claims are
        extrapolated hypotheses.  The critical next step is training
        the tiny variant and the ablated baseline on ImageNet-1K for
        300 epochs (supervised protocol in
        \texttt{TRAINING.md}).
  \item \textbf{No real-dataset results.}  Synthetic micro-benchmarks
        (noise entropy, linear probe on Gaussian blobs, translation
        sweep) are weak proxies for actual ImageNet accuracy.  They
        verify structural correctness but do not predict trained
        performance.
  \item \textbf{FCMAE compatibility untested.}  LE-GRN's local mean
        subtraction may interact poorly with the zeroed-out regions
        in masked autoencoding~\cite{he2022mae}.  A dedicated FCMAE
        pretraining experiment is required.
  \item \textbf{Blocking unknown: uniform downscale baseline.}  The
        most important falsification test---whether a simple uniform
        downscale of ConvNeXt V2-T (channels $[80,160,320,640]$,
        uniform $\times4$) matches V2-E at the same parameter
        budget---has not been trained.
  \item \textbf{MRF-DW wall-clock impact.}  The 4--8\% FLOPs increase
        from the second depthwise branch may cause a larger wall-clock
        slowdown than the FLOPs fraction suggests, since depthwise
        convolutions are memory-bandwidth-bound.
  \item \textbf{Stage~3 parameter concentration.}  With 61.5\% of
        parameters in Stage~3, any underperformance in this stage
        disproportionately affects the overall model.
\end{itemize}

\paragraph{Design-space implications.}
The parameter verification confirms three engineering insights relevant
to efficient ConvNet design: (i) learned per-channel spatial mixing
adds negligible overhead and can be cleanly toggled; (ii) local
contrast enhancement before global normalization is a zero-cost
modification that changes the feature competition dynamics (verified:
LE-GRN outputs differ from GRN on textured synthetic inputs); and
(iii) the uniform $\times4$ expansion inherited from Swin-T is wasteful
at small scales---saving $\sim$1M parameters through stage-adaptive
ratios does not require architectural complexity.

% ============================================================================
\section{Conclusion}
\label{sec:conclusion}

We have presented ConvNeXt V2-E, a reference architecture and
validation harness for efficient ConvNet design in the sub-20M
parameter regime.  The architecture introduces three innovations
targeting specific suboptimalities in ConvNeXt V2 at small scale:
Multi-Receptive Field Depthwise Convolution ($\sim$150K overhead),
Local-Enhanced Global Response Normalization (zero parameter cost),
and stage-adaptive expansion ratios (saving $\sim$0.95M parameters).
The full tiny variant achieves 19.67M parameters and $\sim$3.8G FLOPs
---a 16\% FLOPs reduction versus ConvNeXt V2-T while reallocating
capacity toward the critical middle stage.

The validation harness---comprising shape tests, gradient-flow checks,
numerical stability tests, nine single-field ablations, synthetic
micro-benchmarks, and profiling infrastructure---provides a
reproducible foundation for follow-up work.  All source code,
configuration dataclasses, and evaluation scripts are released under an
open-source license.  \textbf{No trained-model evaluation on
ImageNet-1K has been performed.}

Concrete follow-up directions include:

\begin{enumerate}[nosep]
  \item \textbf{Supervised training on ImageNet-1K} (300~epochs).
        This is the single most important next step: training the tiny
        variant and the ablated baseline to measure each innovation's
        accuracy contribution and to resolve the blocking unknown of
        whether ConvNeXt V2-E justifies its architectural complexity
        over a simple uniform downscale of ConvNeXt V2-T.
  \item \textbf{FCMAE self-supervised pretraining} (400~epochs
        pretraining + 100~epochs fine-tuning).  This would test LE-GRN
        compatibility with masked autoencoding and allow a
        direct comparison with the published ConvNeXt V2-T FCMAE
        result (82.8\% top-1).
  \item \textbf{Downstream dense prediction evaluations} on COCO
        (object detection, instance segmentation) and ADE20K (semantic
        segmentation) to verify that the multi-scale spatial features
        from MRF-DW translate to improvements in dense prediction
        tasks.
\end{enumerate}

By releasing a clean, typed, and thoroughly validated reference
implementation, we aim to accelerate research into efficient ConvNet
architectures at the sub-20M design point---a regime that remains
underexplored in the modern ConvNet literature.

% ============================================================================
% Acknowledgments (omitted in anonymous submission)
% ============================================================================
% \begin{ack}
% This work was supported by ...
% \end{ack}

% ============================================================================
\section*{References}

{\small
\begin{thebibliography}{99}

\bibitem{chen2017deeplab}
L.-C.~Chen, G.~Papandreou, I.~Kokkinos, K.~Murphy, and A.~L. Yuille.
DeepLab: Semantic image segmentation with deep convolutional nets,
atrous convolution, and fully connected CRFs.
\textit{IEEE TPAMI}, 40(4):834--848, 2017.

\bibitem{dosovitskiy2020vit}
A.~Dosovitskiy, L.~Beyer, A.~Kolesnikov, D.~Weissenborn, X.~Zhai,
T.~Unterthiner, M.~Dehghani, M.~Minderer, G.~Heigold, S.~Gelly,
J.~Uszkoreit, and N.~Houlsby.
An image is worth $16\times16$ words: Transformers for image
recognition at scale.
In \textit{ICLR}, 2021.

\bibitem{guo2022visual}
M.-H.~Guo, C.-Z.~Lu, Z.-N.~Liu, M.-M.~Cheng, and S.-M.~Hu.
Visual attention network.
\textit{Computational Visual Media}, 9(4):733--752, 2023.

\bibitem{he2022mae}
K.~He, X.~Chen, S.~Xie, Y.~Li, P.~Doll\'{a}r, and R.~Girshick.
Masked autoencoders are scalable vision learners.
In \textit{CVPR}, 2022.

\bibitem{howard2019mobilenetv3}
A.~Howard, M.~Sandler, G.~Chu, L.-C.~Chen, B.~Chen, M.~Tan,
W.~Wang, Y.~Zhu, R.~Pang, V.~Vasudevan, Q.~V. Le, and H.~Adam.
Searching for MobileNetV3.
In \textit{ICCV}, 2019.

\bibitem{hu2018squeeze}
J.~Hu, L.~Shen, and G.~Sun.
Squeeze-and-excitation networks.
In \textit{CVPR}, 2018.

\bibitem{krizhevsky2012alexnet}
A.~Krizhevsky, I.~Sutskever, and G.~E. Hinton.
ImageNet classification with deep convolutional neural networks.
In \textit{NeurIPS}, 2012.

\bibitem{liu2021swin}
Z.~Liu, Y.~Lin, Y.~Cao, H.~Hu, Y.~Wei, Z.~Zhang, S.~Lin, and
B.~Guo.
Swin Transformer: Hierarchical vision transformer using shifted
windows.
In \textit{ICCV}, 2021.

\bibitem{liu2022convnext}
Z.~Liu, H.~Mao, C.-Y.~Wu, C.~Feichtenhofer, T.~Darrell, and S.~Xie.
A ConvNet for the 2020s.
In \textit{CVPR}, 2022.

\bibitem{radosavovic2020regnet}
I.~Radosavovic, R.~P. Kosaraju, R.~Girshick, K.~He, and P.~Doll\'{a}r.
Designing network design spaces.
In \textit{CVPR}, 2020.

\bibitem{szegedy2015inception}
C.~Szegedy, W.~Liu, Y.~Jia, P.~Sermanet, S.~Reed, D.~Anguelov,
D.~Erhan, V.~Vanhoucke, and A.~Rabinovich.
Going deeper with convolutions.
In \textit{CVPR}, 2015.

\bibitem{tan2019efficientnet}
M.~Tan and Q.~V. Le.
EfficientNet: Rethinking model scaling for convolutional neural
networks.
In \textit{ICML}, 2019.

\bibitem{tan2019mixconv}
M.~Tan and Q.~V. Le.
MixConv: Mixed depthwise convolutional kernels.
In \textit{BMVC}, 2019.

\bibitem{ulyanov2016instance}
D.~Ulyanov, A.~Vedaldi, and V.~Lempitsky.
Instance normalization: The missing ingredient for fast stylization.
arXiv:1607.08022, 2016.

\bibitem{woo2023convnextv2}
S.~Woo, S.~Debnath, R.~Hu, X.~Chen, Z.~Liu, I.~S. Kweon, and S.~Xie.
ConvNeXt V2: Co-designing and scaling ConvNets with masked
autoencoders.
In \textit{CVPR}, 2023.

\bibitem{yang2022focal}
J.~Yang, C.~Li, P.~Zhang, X.~Dai, B.~Xiao, L.~Yuan, and J.~Gao.
Focal modulation networks.
In \textit{NeurIPS}, 2022.

\end{thebibliography}
}

% ============================================================================
\appendix
\section{NeurIPS Paper Checklist}

\paragraph{1. Claims.}
Do the main claims made in the abstract and introduction accurately
reflect the paper's contributions and scope?
\textbf{Yes.}  The paper claims only the design artifact (architecture
+ validation harness) and explicitly states that no trained-model
evaluation on real datasets has been performed.

\paragraph{2. Experiments.}
Does the paper provide sufficient information to reproduce the
experiments?
\textbf{Yes.}  Every test, benchmark, and ablation is associated with a
specific file and command in the validation harness.  The
\texttt{ModelConfig} dataclass and variant constructors contain all
hyperparameters without magic numbers.

\paragraph{3. Theory.}
Does the paper state the theoretical basis for its claims?
\textbf{Yes.}  Parameter formulas are derived per-block
(Eq.~\ref{eq:block_params}), per-innovation
(Eq.~\ref{eq:mrf_overhead}, \ref{eq:adapt_savings}), and per-stage.
Falsification conditions are specified for each architectural
innovation.

\paragraph{4. Datasets.}
Does the paper use any datasets?
\textbf{No.}  The validation harness uses synthetic random inputs and
Gaussian blob data.  ImageNet-1K is referenced as the target for
follow-up training but is not used in this work.

\paragraph{5. Code.}
Is the code publicly available?
\textbf{Yes.}  The reference implementation, configuration dataclasses,
test suite, ablation runner, benchmarks, and profiling scripts are
released as part of the paper artifact bundle.

\paragraph{6. Broader impacts.}
What are the broader impacts of the work?
The paper presents an architectural design and validation methodology
for efficient ConvNets.  Potential positive impacts include enabling
deployment of modern ConvNet backbones on resource-constrained devices.
Potential negative impacts are not foreseen beyond the general
environmental considerations of training deep neural networks.

\end{document}
